Este capítulo aborda la fase de análisis y diseño del proyecto, estableciendo los cimientos técnicos y lógicos sobre los cuales se ha construido la plataforma SocratiCode. En primer lugar, se detalla la metodología de trabajo adoptada para la gestión y desarrollo del proyecto (\autoref{SEC:AD_METODOLOGIA}), seguida de una descripción general del propósito y alcance del sistema (\autoref{SEC:AD_DESCRIPCION}). A continuación, se definen las necesidades del software mediante la especificación de sus requisitos funcionales y no funcionales (\autoref{SEC:AD_REQUISITOS}), los cuales se materializan posteriormente a través de los diferentes casos de uso (\autoref{SEC:AD_CASOSUSO}). Finalmente, para consolidar la visión técnica de la aplicación, se presenta el modelo de datos empleado (\autoref{SEC:AD_MODELO}) junto con el diagrama de componentes y arquitectura a alto nivel (\autoref{SEC:AD_ARQ}), cerrando el capítulo con unas breves conclusiones sobre esta fase (\autoref{SEC:AD_CONCLUSIONES}).
